% Contenidos del capítulo.
% Las secciones presentadas son orientativas y no representan
% necesariamente la organización que debe tener este capítulo.

\section{Introducción}
La confluencia de tres vectores tecnológicos la maduración del \gls{agenteia}, la viabilidad de la inferencia local sobre \gls{hardware} de consumo y la creciente exigencia social y normativa de privacidad en el tratamiento de datos está redefiniendo el paradigma del desarrollo de sistemas inteligentes \cite{xi2023rise}. Hasta hace pocos años, la \gls{ia} avanzada exigía infraestructuras de computación en la nube, con los costes de latencia, dependencia del proveedor y exposición de datos que ello conlleva. Hoy, los \glspl{llm} cuantizados, las arquitecturas de razonamiento multitarea y los protocolos abiertos de herramientas permiten concebir asistentes autónomos que operan íntegramente en el dispositivo del usuario, bajo el paradigma \gls{localfirst}. Este cambio no es meramente incremental, sino que traslada el centro de gravedad del ecosistema de \gls{ia} desde unos pocos proveedores centralizados hacia una red distribuida de agentes personales gobernados por sus propios usuarios, con profundas implicaciones tanto para la privacidad individual como para la soberanía tecnológica de organizaciones y administraciones públicas.

El presente \gls{tfm} se sitúa precisamente en la intersección de estas tres líneas de fuerza. Su propósito es explorar tanto en el plano conceptual como en el de la ingeniería cómo combinar agentes con capacidad de razonamiento encadenado, recuperación semántica de información y ejecución de herramientas externas, todo ello desplegado en un entorno de ejecución local, reproducible y transparente. En las secciones siguientes se desglosan los fundamentos teóricos que sustentan esta propuesta, se exponen las motivaciones que la impulsan y se enuncian los objetivos que guían el desarrollo del proyecto LocalMind-AI.

\subsection{Agentes de IA y Arquitecturas de Razonamiento}
\label{subsec:agentes-ia}

El punto de partida conceptual del proyecto es la noción de \gls{agenteia}, un sistema que percibe su entorno a través de sensores ya sean físicos, como cámaras o micrófonos, o digitales, como \glspl{api} y flujos de datos, razona sobre el estado de dicho entorno mediante modelos internos y ejecuta acciones orientadas a la consecución de uno o varios objetivos predefinidos o emergentes \parencite[p.~2]{xi2023rise}. Este bucle percepción-razonamiento-acción, ilustrado en la Figura \ref{fig:agenteEntornoBucle}, constituye el núcleo funcional de cualquier arquitectura de agente autónomo y establece una separación conceptual clara entre el flujo de información sensorial entrante, el procesamiento deliberativo interno y la actuación sobre el mundo. La calidad del razonamiento interno la capacidad del agente para interpretar correctamente las observaciones, anticipar consecuencias y seleccionar la acción adecuada determina, en última instancia, su eficacia y su robustez frente a entornos dinámicos.

\begin{figure}[h!tp]
    \centering
    \includegraphics[width=0.65\textwidth]{agenteEntornoBucle.png}
    \caption{Bucle percepción-razonamiento-acción de un agente inteligente.}
    \label{fig:agenteEntornoBucle}
\end{figure}

En el plano del razonamiento, la técnica de \gls{cot} \cite{wei2022chain} supuso un punto de inflexión al demostrar que los \glspl{llm} mejoran sustancialmente su rendimiento en tareas que requieren múltiples pasos lógicos cuando se les induce mediante ejemplos en el \gls{prompt} o instrucciones explícitas a descomponer el problema en etapas intermedias y a verbalizar el proceso antes de emitir una respuesta. Este hallazgo reveló que la generación paso a paso no solo incrementa la precisión, sino que también dota de trazabilidad a las decisiones del modelo, facilitando la auditoría y la depuración del comportamiento del agente.

Sobre esta base, la arquitectura \gls{react} \cite{yao2022react} introdujo un salto cualitativo al entrelazar explícitamente los pasos de razonamiento interno con la invocación de herramientas externas. En el bucle \gls{react}, representado en la Figura \ref{fig:bucleReAct}, el agente alterna fases de pensamiento durante las cuales analiza la situación, formula hipótesis y planifica con fases de acción en las que consulta una \gls{api}, ejecuta código, recupera documentos o interactúa con una base de datos, realimentándose con los resultados observados hasta alcanzar el objetivo. Esta alternancia confiere al agente una capacidad de adaptación dinámica que los modelos puramente generativos no poseen, ya que el flujo de razonamiento se corrige y reorienta en tiempo real a partir de la información obtenida del entorno.

\begin{figure}[h!tp]
    \centering
    \includegraphics[width=0.65\textwidth]{bucleReAct.png}
    \caption{Bucle \gls{react} con razonamiento y acción entrelazados con el entorno.}
    \label{fig:bucleReAct}
\end{figure}

Más allá de la alternancia razonamiento-acción, la investigación reciente ha incorporado mecanismos de \gls{reflexion} \cite{shinn2023reflexion}, que permiten al \gls{agenteia} evaluar su propio rendimiento tras cada ciclo de ejecución, detectar errores o ineficiencias y ajustar su estrategia en iteraciones sucesivas sin intervención humana. Esta capacidad metacognitiva, aunque todavía incipiente en los sistemas actuales, apunta hacia \glspl{agenteia} capaces de aprendizaje continuo en entornos abiertos sin requerir reentrenamiento supervisado, un avance que tendría implicaciones profundas para la autonomía de los asistentes personales. De forma complementaria, las técnicas de \gls{planeacion} entre las que destaca \textit{Tree of Thoughts} \cite{yao2023tree} dotan al \gls{agenteia} de la capacidad de explorar múltiples líneas de razonamiento en paralelo, evaluar su viabilidad mediante heurísticas de búsqueda y seleccionar la trayectoria más prometedora antes de comprometer acciones irreversibles. La combinación de estos paradigmas encadenamiento de pensamiento, actuación entrelazada, autorreflexión y \gls{planeacion} configura un espacio de diseño rico en el que cada decisión arquitectónica afecta al equilibrio entre precisión, latencia y consumo de recursos computacionales.

En síntesis, las arquitecturas de razonamiento descritas se organizan en torno a cuatro paradigmas complementarios.

\begin{itemize}
    \item \textbf{\gls{cot}.} Descomposición paso a paso del problema en etapas intermedias, que incrementa \cite{wei2022chain} la precisión y dota de trazabilidad a las decisiones del modelo.
    \item \textbf{\gls{react}.} Entrelazado de pasos de razonamiento interno con herramientas externas, realimentándose en tiempo real con los resultados observados \cite{yao2022react}.
    \item \textbf{\gls{reflexion}.} Autoevaluación del desempeño tras cada ciclo de ejecución para detectar errores y ajustar la estrategia sin intervención humana \cite{shinn2023reflexion}.
    \item \textbf{\gls{planeacion}.} Exploración de múltiples líneas de razonamiento en paralelo con selección de la trayectoria más prometedora \cite{yao2023tree}.
\end{itemize}

El puente entre el razonamiento simbólico y la ejecución concreta se materializa mediante el mecanismo de \gls{funciontool}, que permite al modelo de lenguaje invocar procedimientos externos parametrizados consultas a bases de datos, llamadas a servicios web, ejecución de \textit{scripts} y utilizar los resultados para enriquecer el contexto de razonamiento. Esta capacidad de \textit{tool use} dota al \gls{agenteia} de un campo de acción que trasciende la generación de texto y lo aproxima al comportamiento de un sistema operativo cognitivo, capaz de orquestar recursos heterogéneos en función de los objetivos que persigue. La eficacia de este mecanismo depende críticamente de la calidad de la definición de cada herramienta su nombre, descripción y esquema de parámetros, ya que el modelo debe inferir a partir de estas señales cuándo y cómo invocarla, sin conocimiento previo de su implementación interna.

En el extremo superior del espectro de complejidad se sitúan los \glspl{sistemaMultiagente} \cite{xi2023rise}, compuestos por múltiples \glspl{agenteia} que colaboran, negocian o compiten entre sí para abordar problemas cuya descomposición natural excede las capacidades de una entidad única. Estos sistemas introducen desafíos adicionales de coordinación, comunicación y resolución de conflictos que los convierten en un campo de investigación activo y en rápida evolución. Aunque el proyecto que aquí se presenta no adopta una arquitectura de \glspl{sistemaMultiagente}, conviene señalar que este enfoque representa la frontera actual de la investigación en agentes autónomos y constituye una dirección natural de evolución para los asistentes locales, especialmente en escenarios donde la especialización de \glspl{agenteia} por dominio puede ofrecer ventajas cualitativas frente a arquitecturas monolíticas.

\subsection{Entornos de Ejecución Local y Privacidad de Datos}
\label{subsec:ejecucion-local}

Durante la última década, el modelo dominante para el despliegue de modelos de \gls{ia} ha sido la computación en la nube, donde los datos del usuario, consultas, documentos e imágenes se transmiten a centros de datos remotos donde residen los modelos, se ejecuta la inferencia sobre \gls{hardware} especializado y se devuelve el resultado. Este esquema ofrece escalabilidad prácticamente ilimitada y acceso a aceleradores de última generación, pero introduce tres fricciones que resultan cada vez más difíciles de ignorar, una dependencia total al proveedor que limita su uso al usuario, una latencia de red que penaliza las aplicaciones interactivas y sobre todo, una cesión de control sobre los datos potencialmente sensibles que resulta incompatible con los estándares contemporáneos de privacidad \cite{kleppmann2019local}. La Figura \ref{fig:localVsCloud} contrapone ambos modelos de forma esquemática.

\begin{figure}[h!tp]
    \centering
    \includegraphics[width=0.65\textwidth]{localVsCloud.png}
    \caption{Comparativa entre inferencia en la nube e inferencia local.}
    \label{fig:localVsCloud}
\end{figure}

Frente a esta realidad, el paradigma \gls{localfirst} propone invertir dicha relación, los datos y la computación priman en el dispositivo del usuario y la conexión a servicios remotos se reserva para funciones auxiliares estrictamente necesarias, como la sincronización entre dispositivos o la colaboración en tiempo real \cite{kleppmann2019local}. La viabilidad técnica de este enfoque ha sido impulsada por dos avances convergentes. Por un lado, las técnicas de cuantización de modelos que reducen la precisión numérica de los pesos desde representaciones de 16 o 32 bits hasta 4 u 8 bits sin una degradación significativa del rendimiento permiten ejecutar \glspl{llm} con miles de millones de parámetros en equipos portátiles con memorias de 16 GB o incluso inferiores. Por otro, la incorporación de unidades de procesamiento neuronal en los procesadores de nueva generación y de manera particular, la aceleración por \gls{gpu} a través de la \gls{api} Metal en los chips Apple Silicon presentes en la gama M1, M2, M3 y M4 o el uso de \glspl{gpu} de Nvidia proporcionan el rendimiento necesario para que la inferencia local sea interactiva, con tasas de generación de tokens comparables a las de servicios en la nube para modelos de tamaño moderado. Diversas herramientas de código abierto han democratizado este acceso al ofrecer interfaces unificadas para la descarga, cuantización y ejecución de decenas de modelos sobre \gls{hardware} convencional \cite{ollama2024}, eliminando la barrera de entrada que suponía la configuración manual de entornos de inferencia.

Las implicaciones de este desplazamiento hacia el extremo del dispositivo trascienden la mera ingeniería de sistemas. Desde el punto de vista normativo, la ejecución local simplifica el cumplimiento de marcos regulatorios como el \gls{rgpd}, la \gls{lopdgdd} y el \gls{aiact} \parencite{ReglamentoUE20162016}, al reducir o eliminar la exposición de datos personales a terceros y al garantizar que el tratamiento se produce bajo el control exclusivo del usuario, sin que la información abandone su dispositivo. La combinación de inferencia local, modelos abiertos y despliegue contenerizado sienta las bases de un ecosistema de asistentes de \gls{ia} que pueden ser auditados, reproducidos y gobernados por sus propios usuarios, en lugar de por proveedores externos cuyos intereses comerciales no siempre coinciden con los del usuario final.

\subsection{Modelos de Lenguaje, Recuperación Aumentada y Herramientas Externas}
\label{subsec:llm-rag}

Si el bucle percepción-razonamiento-acción constituye el esqueleto del agente, el \gls{llm} representa su núcleo cognitivo. Los modelos de lenguaje de gran escala han demostrado una capacidad notable para comprender instrucciones en lenguaje natural, generar texto coherente en múltiples idiomas y, cuando se combinan con las arquitecturas de razonamiento descritas anteriormente, descomponer problemas complejos en secuencias de acciones ejecutables \cite{wei2022chain}. No obstante, el conocimiento de un \gls{llm} está inevitablemente acotado por los datos con los que fue entrenado, lo que introduce dos limitaciones críticas, como son la imposibilidad de acceder a información posterior a la fecha de corte del corpus de entrenamiento y la propensión a generar contenido plausible pero factualmente incorrecto fenómeno conocido como alucinación, especialmente cuando el modelo se enfrenta a preguntas que requieren datos precisos o actualizados.

La arquitectura de \gls{rag} \cite{lewis2020retrieval} aborda ambas limitaciones mediante la introducción de un módulo de recuperación de información externa que se ejecuta previamente a la generación de la respuesta. El proceso, representado de forma esquemática en la Figura \ref{fig:pipelineRag}, se articula en dos fases diferenciadas y complementarias. En la fase de indexado que se realiza típicamente de manera asíncrona y por lotes, un conjunto de documentos de referencia se fragmenta en unidades textuales manejables, cada fragmento se transforma en un \gls{embeddings} mediante un modelo especializado de codificación semántica, y el conjunto de vectores resultante se almacena en una base de datos vectorial optimizada para búsquedas por similitud. En la fase de consulta, la pregunta del usuario se convierte a su vez en un vector de consulta utilizando el mismo modelo de codificación, se recuperan de la base de datos vectorial los fragmentos más relevantes por proximidad semántica típicamente los $k$ vectores más cercanos y dichos fragmentos se inyectan como contexto adicional en el \gls{prompt} que se envía al \gls{llm}. De este modo, el modelo genera su respuesta fundamentándose en información recuperada de fuentes externas y verificables, lo que incrementa significativamente tanto la precisión factual como la trazabilidad de las respuestas, ya que el usuario puede inspeccionar los fragmentos que sirvieron de base para la generación. La calidad del sistema \gls{rag} depende de tres factores interrelacionados, que corresponden a la estrategia de fragmentación de los documentos que debe equilibrar la granularidad semántica con la coherencia textual, el modelo de \glspl{embeddings} empleado cuya capacidad de capturar relaciones semánticas determina la pertinencia de los fragmentos recuperados y el criterio de selección de los $k$ vecinos más cercanos, que regula el equilibrio entre exhaustividad y concisión del contexto inyectado en el \gls{prompt}.

En síntesis, el pipeline \gls{rag} se articula en dos fases complementarias

\begin{itemize}
    \item \textbf{Fase de indexado} los documentos de referencia se fragmentan, se transforman en \glspl{embeddings} mediante un modelo de codificación semántica y se almacenan en una base de datos vectorial optimizada para búsquedas por similitud.
    \item \textbf{Fase de consulta} la pregunta del usuario se convierte en un vector, se recuperan los $k$ fragmentos más relevantes por proximidad semántica y se inyectan como contexto en el \gls{prompt} del \gls{llm} para fundamentar la respuesta \cite{lewis2020retrieval}.
\end{itemize}

\begin{figure}[h!tp]
    \centering
    \includegraphics[width=0.65\textwidth]{pipelineRag.png}
    \caption{Pipeline RAG: indexado, recuperación y generación aumentada.}
    \label{fig:pipelineRag}
\end{figure}

La extensión lógica del paradigma \gls{rag} consiste en otorgar al modelo no solo acceso a documentos, sino a herramientas computacionales completas que le permitan actuar sobre el mundo digital. En este contexto surge el \gls{mcp} \cite{anthropic2024mcp}, un protocolo abierto diseñado para estandarizar la forma en que los modelos de lenguaje descubren, invocan y reciben resultados de herramientas externas. \Gls{mcp} establece una separación explícita entre el proveedor del modelo quien suministra la capacidad de razonamiento y el proveedor de las herramientas quien expone funcionalidades concretas como ejecución de código, generación de documentos o acceso a sistemas de archivos, definiendo un contrato de comunicación basado en mensajes estructurados que cualquier servidor compatible puede implementar. Esta desacoplación permite que el agente integre dinámicamente nuevas capacidades en su flujo de razonamiento sin modificar su núcleo, favoreciendo la extensibilidad y la modularidad del sistema. El protocolo define, además, mecanismos de negociación de capacidades y transporte de contexto que permiten al agente mantener un modelo coherente de las herramientas disponibles y de su estado en cada momento, lo que resulta esencial para la toma de decisiones informada en flujos de razonamiento de múltiples pasos.

La conjunción de \gls{rag} como mecanismo de fundamentación factual y \gls{mcp} como capa estandarizada de herramientas configura una arquitectura en la que el \gls{llm} deja de ser un mero generador de texto para convertirse en un orquestador cognitivo que razona, recupera información pertinente, decide qué herramienta invocar en cada paso e interpreta los resultados obtenidos, todo ello dentro de un bucle coherente, trazable y potencialmente verificable por el usuario.

\section{Motivación}
\label{sec:motivacion}

Los avances teóricos expuestos en la sección anterior representan la base tecnológica sobre la que es posible construir hoy asistentes de \gls{ia} que operen de forma autónoma, privada y reproducible sobre \gls{hardware} de consumo. Sin embargo, la integración efectiva de estos componentes (un núcleo de razonamiento \gls{cot}-\gls{rag}, un conjunto de herramientas externas gobernadas por \gls{mcp}, una interfaz de usuario accesible localmente desde la web y un despliegue local contenerizado que garantice la reproducibilidad) sigue siendo un desafío de ingeniería que exige tomar decisiones de diseño concretas, implementarlas con rigor y validarlas empíricamente frente a escenarios de uso realistas.

El proyecto LocalMind-AI nace precisamente de la convicción de que un asistente personal de \gls{ia} no debería depender de servicios en la nube para ofrecer respuestas de calidad, ni obligar al usuario a ceder el control sobre sus datos a cambio de funcionalidad. La elección de \gls{nanobot} como \textit{framework} de orquestación del \gls{agenteia} responde a su perfil ultraligero diseñado para entornos con recursos limitados y a su compatibilidad nativa con los ecosistemas de \gls{ollama} y \gls{omlx}, lo que permite que el \gls{agenteia} y el modelo de lenguaje coexistan en el mismo dispositivo sin tensionar los recursos del sistema. La interfaz web, por su parte, se construye sobre \gls{reactnative} \cite{meta2024reactnative} y \gls{expo} \cite{expo2024} con el propósito de garantizar un acceso unificado desde cualquier navegador y facilitar la iteración rápida durante las fases de prototipado y validación. En última instancia, este \gls{tfm} persigue demostrar la viabilidad y efectividad de una arquitectura de \gls{agenteia} local, privada y modular en \gls{hardware} de consumo, ofreciendo una alternativa robusta a las soluciones comerciales centralizadas en la nube y garantizando que tanto el flujo de decisión como la gestión de datos permanezcan bajo el control absoluto y auditable del usuario.

\section{Objetivos}
\label{sec:objetivos}

El objetivo principal de este \gls{tfm} es el diseño, la implementación y la validación experimental de un asistente de \gls{ia} multitarea que integre razonamiento \gls{cot}-\gls{rag} iterativo, memoria persistente a largo plazo, invocación de herramientas externas a través del protocolo \gls{mcp} y un despliegue reproducible mediante \gls{docker} \gls{dockercompose} \cite{docker2024compose}. Este despliegue separa el \gls{agenteia} y el cliente web en dos contenedores, mientras que los modelos \gls{llm} locales gestionados por \gls{ollama} u \gls{omlx} se ejecutan en el anfitrión para aprovechar su aceleración. La solución comprende una interfaz web desarrollada con \gls{reactnative} y \gls{expo}, exportada como contenido estático y servida por \gls{nginx}, y está orientada a usuarios con recursos computacionales limitados como un ordenador portátil de gama media o un dispositivo de consumo que requieran asistencia automatizada para tareas diversas preservando en todo momento la privacidad y la soberanía sobre sus datos, sin depender de conexiones de red de alta velocidad ni de suscripciones a servicios externos de \gls{ia} \cite{nginx2026}.

\subsection{Objetivos Específicos y Retos Técnicos}
\label{subsec:objetivos-especificos}

Para alcanzar el objetivo general, se abordan los siguientes retos técnicos mediante una metodología estructurada que permite la verificación incremental de cada componente.

\begin{enumerate}
    \item \textbf{Implementación de un flujo de razonamiento CoT-RAG con memoria persistente sobre modelos locales.} Integrar el bucle iterativo \gls{cot}-\gls{rag} y la memoria de largo plazo dentro de un \gls{agenteia} orquestado por \gls{nanobot}, utilizando \gls{ollama} u \gls{omlx} como motores de inferencia para modelos cuantizados, \gls{chromadb} como almacén vectorial para la indexación y recuperación semántica de documentos, y \gls{engram} como sistema de memoria persistente para el almacenamiento de observaciones y notas del \gls{agenteia}. El flujo debe ser capaz de alternar automáticamente entre pasos de razonamiento interno, recuperación de información y consulta de memoria, inyectando el contexto relevante en el \gls{prompt} del modelo de forma transparente para el usuario.

    \item \textbf{Integración de herramientas externas mediante el Model Context Protocol.} Diseñar e implementar un servidor \gls{mcp} personalizado que exponga capacidades de generación de documentos \gls{pdf} e interacción con la base de conocimiento local, permitiendo al \gls{agenteia} invocarlas de manera autónoma cuando su flujo de razonamiento lo requiera.

    \item \textbf{Construcción de una interfaz web con React Native y Expo.} Desarrollar una aplicación web que ofrezca un chat conversacional con visualización del flujo de razonamiento incluyendo los pasos intermedios de \gls{cot}, gestión de tareas pendientes y completadas, y configuración dinámica de los parámetros del asistente (modelo, dirección del servidor \gls{llm} y herramientas activas), garantizando accesibilidad desde cualquier navegador moderno mediante una base de código unificada.

    \item \textbf{Despliegue reproducible del sistema completo mediante Docker Compose.} Definir un entorno de ejecución con dos servicios diferenciados, uno para el \gls{agenteia} \gls{nanobot} y otro para el cliente web. El primero se construye mediante etapas separadas y se ejecuta con el usuario no privilegiado \texttt{nanobot}, mientras que el segundo exporta la aplicación \gls{expo} en una etapa de construcción con \gls{nodejs} y \gls{pnpm} y sirve los activos resultantes mediante \gls{nginx} \cite{docker2026multistage,nginx2026}. Los servidores \gls{mcp} y \gls{engram} se inician como subprocesos del agente, y \gls{ollama} u \gls{omlx} permanecen en el sistema anfitrión para aprovechar la aceleración por \gls{gpu}. De este modo, cualquier usuario puede reproducir los componentes contenerizados en su propio \gls{hardware} con un único comando de \gls{dockercompose}.

    \item \textbf{Validación funcional y de rendimiento del asistente.} Diseñar y ejecutar un plan de pruebas sistemático que cubra la corrección del flujo \gls{cot}-\gls{rag} verificando que las respuestas incorporan información de los documentos recuperados, la latencia de extremo a extremo, la precisión de la recuperación semántica frente a un conjunto de referencia y la estabilidad del sistema bajo cargas de trabajo representativas, documentando tanto los resultados favorables como las limitaciones identificadas.
\end{enumerate}

\subsection{Metodología de Desarrollo}
\label{subsec:metodologia}

La ejecución del proyecto se adhiere a una \gls{iterativa} \cite{larmanIterativeIncrementalDevelopments2003}, inspirado en los principios de la \gls{metagil} \cite{ilievaAnalysesAgileMethodology2004}, que organiza el trabajo en ciclos progresivos de refinamiento y entrega incremental de valor. Esta elección metodológica se justifica por la naturaleza exploratoria del proyecto, en el que tanto la selección de tecnologías modelos de lenguaje, \textit{frameworks} de orquestación, protocolos de herramientas como la definición de la arquitectura evolucionan a medida que se adquiere conocimiento empírico sobre el comportamiento de los componentes y su interacción. Además, el enfoque iterativo permite incorporar hallazgos de cada ciclo de validación en la iteración siguiente, reduciendo el riesgo de descubrir problemas estructurales en fases tardías del desarrollo.

El desarrollo se estructura en cuatro fases principales, cada una de las cuales produce artefactos verificables que alimentan la fase siguiente y permiten la realimentación temprana sobre las decisiones adoptadas. La fase de investigación comprende la revisión sistemática del estado del arte en \glspl{agenteia}, arquitecturas de razonamiento \gls{cot}, \gls{react}, \gls{reflexion}, sistemas \gls{rag}, el protocolo \gls{mcp} y \textit{frameworks} de orquestación, así como la comparativa de modelos \gls{llm} cuantizados y la medición de su rendimiento en \gls{hardware} local con y sin aceleración por \gls{gpu}. La fase de arquitectura aborda la selección definitiva de tecnologías y el diseño detallado de los componentes del sistema, como el \gls{agenteia} basado en \gls{nanobot} con su \textit{skill} \gls{cot}-\gls{rag}, los subprocesos \gls{mcp} y \gls{engram}, la base de datos vectorial \gls{chromadb} y la aplicación web en \gls{reactnative} con \gls{expo}, incluyendo los diagramas de componentes, secuencia y despliegue que documentan las decisiones tomadas. La fase de implementación materializa cada componente de forma incremental e integra en \gls{dockercompose} dos servicios para el agente y el cliente web, manteniendo la inferencia acelerada en el anfitrión. Por último, la fase de validación somete el sistema integrado a pruebas funcionales y de rendimiento, documenta los resultados obtenidos y extrae conclusiones de cara a futuras iteraciones o líneas de trabajo derivadas.

\section{Organización de la Memoria}
\label{sec:organizacion}

La presente memoria se estructura en siete capítulos que recorren de forma progresiva el fundamento teórico, el análisis, el diseño, la implementación y la evaluación del proyecto, facilitando una lectura lineal que va de lo general a lo particular.

\begin{itemize}
    \item \textbf{Capítulo 1.} Presenta el contexto tecnológico en el que se inscribe LocalMind-AI (\glspl{agenteia}, razonamiento encadenado, ejecución local y herramientas externas), establece la motivación que impulsa el trabajo, define los objetivos generales y específicos, y describe la organización del resto de la memoria.

    \item \textbf{Capítulo 2.} Expone el estado del arte en las áreas relevantes para el proyecto, abarcando \glspl{agenteia} y arquitecturas de razonamiento \gls{cot}, \gls{react} y \gls{reflexion}, sistemas \gls{rag}, el protocolo \gls{mcp}, servidores de inferencia local como \gls{ollama} y \gls{omlx}, memoria persistente con \gls{engram}, \glspl{llm} locales y cuantización, \textit{frameworks} de orquestación como \gls{nanobot} y tecnologías de desarrollo multiplataforma con soporte web como \gls{reactnative} y \gls{expo}, estableciendo una comparativa crítica entre las alternativas existentes.

    \item \textbf{Capítulo 3.} Especifica los \glspl{requisitosfuncionales} y \glspl{requisitosnofuncionales} del sistema, detalla la planificación temporal y la estimación de costes del proyecto, identifica los riesgos técnicos y de gestión, y analiza la viabilidad técnica, económica y legal del asistente propuesto.

    \item \textbf{Capítulo 4.} Desarrolla el análisis del sistema mediante la identificación de actores y partes interesadas, la definición de \glspl{casosdeuso} que capturan las interacciones esperadas y una primera aproximación al modelo de dominio a través de un \gls{diagraclases} que refleja las entidades y relaciones fundamentales.

    \item \textbf{Capítulo 5.} Describe el diseño arquitectónico del sistema completo, incluyendo la descomposición en componentes del \gls{agenteia}, los subprocesos \gls{mcp} y \gls{engram}, la base de datos vectorial, la aplicación web, los diagramas de comunicación y secuencia, el modelo de datos, el \gls{diagradespliegue} que ilustra la distribución física de los contenedores y las decisiones de diseño razonadas para cada subsistema.

    \item \textbf{Capítulo 6.} Documenta la implementación detallada del \gls{agenteia} \gls{nanobot} y su \textit{skill} \gls{cot}-\gls{rag}, los servidores \gls{mcp} ejecutados como subprocesos, la memoria semántica con \gls{engram}, la base de datos vectorial \gls{chromadb}, la interfaz web \gls{reactnative} con \gls{expo} servida por \gls{nginx} y la orquestación de los servicios contenerizados mediante \gls{docker} \gls{dockercompose}, manteniendo la inferencia en el anfitrión. Asimismo, presenta los resultados de las pruebas funcionales y de rendimiento realizadas.

    \item \textbf{Capítulo 7.} Recoge las conclusiones del trabajo, contrasta los objetivos alcanzados con los planteados inicialmente, revisa los costes finales frente a la estimación inicial y propone líneas de trabajo futuro para la evolución del asistente, incluyendo posibles extensiones hacia \glspl{sistemaMultiagente} y nuevos dominios de aplicación.
\end{itemize}
